% GENERATED FILE. Edit canonical lessons, then run npm run generate:books.
% canonical-source-hash: fnv1a64:fc9c5a70b5bccdbc
% canonical-lessons: ZH-C451-chisu, ZH-C451-tese, ZH-W451-se-colour, ZH-C451-chuse, ZH-C451-sucai, ZH-C451-sheng

\chapter{Specialities}
\label{ch:zh-specialities}
\hlchaptermodality{451}

\begin{chapteropening}
\textbf{By the end of this chapter:} \emph{I can say I eat no meat, talk about a place's special character, say something is outstanding, ask for meatless dishes and take food that is left over.}

Complete the last lesson of chapter 451: I can say I eat no meat, talk about a place's special character, say something is outstanding, ask for meatless dishes and take food that is left over.
\end{chapteropening}

\section[chīsù]{chīsù — to eat no meat; to be vegetarian}

\label{lesson:ZH-C451-chisu}



\begin{quote}
Before the new one: say the Mandarin for authentic, then the Mandarin for a restaurant's signature dish.
\end{quote}

\subsection*{You'll want to know: chīsù}
\emph{chīsù} — \textquotedblleft{}to eat no meat; to be vegetarian\textquotedblright{}.

\begin{quote}
\emph{wǒ chīsù bù chī ròu} — I'm vegetarian: I don't eat meat.
\end{quote}

Its characters come later; for now it is a word for the ear.

\begin{cousinweb}
\emph{chīsù} is built from \zh{吃} 'to eat' and \emph{sù} 'plain'. For the first the dictionary has \emph{khek}, glossed 'to eat', under another way of writing it (there \zh{吃} itself is \emph{kj+t}, 'to stammer'), though \emph{khek} would regularly give qì; the second is inherited from Middle Chinese \emph{suH} (Guangyun, in Baxter's transcription), glossed as plain white silk, and the sense 'meatless' is not given there.
\end{cousinweb}

\begin{grammarlens}[title={What you've built}]
One more word for getting around.
\end{grammarlens}

\subsection*{Your turn}
\begin{itemize}
\raggedright
  \item \emph{Say it:} \emph{chīsù}
  \item \emph{Say it:} \emph{chīsù}, once more
  \item \emph{Say it:} \emph{wǒ chīsù bù chī ròu} — I'm vegetarian: I don't eat meat.
  \item \emph{Recall:} say the Mandarin for authentic, then the Mandarin for a restaurant's signature dish, then say \emph{chīsù} again
\end{itemize}

\begin{tcolorbox}[breakable,colback=teal!4,colframe=teal!35!black,title={Before you move on}]
What does chīsù mean? (\textquotedblleft{}To eat no meat; to be vegetarian\textquotedblright{}.) Say it once more.
\end{tcolorbox}
\section[\texorpdfstring{\zh{特色} (tèsè)}{特色 (tèsè)}]{\zh{特色} (tèsè) — special character; speciality}

\label{lesson:ZH-C451-tese}



\begin{quote}
Before the new one: say the Mandarin for a restaurant's signature dish, then the Mandarin for to eat no meat.
\end{quote}

\subsection*{You'll want to know: \zh{特色}}
\textbf{\zh{特色}} — \emph{tèsè} — \textquotedblleft{}special character; speciality\textquotedblright{}.

\begin{quote}
\emph{zhè jiā fànguǎn de cài hěn yǒu tèsè} — The food at this restaurant has a character all its own.
\end{quote}

The new character, in pieces: \zh{色} — six strokes in two pieces: on top a short stroke falling to the left and a small bar that turns down to the left; under them a bar that turns down, a short vertical, a bar, and a long stroke that runs down the left, along the foot and hooks up. The next lesson puts it in your hand.

\begin{cousinweb}
\textbf{\zh{特色}} (\emph{tèsè}) is built from \zh{特} 'single, special' and \zh{色} 'a look', both inherited from Middle Chinese, \emph{dok} and \emph{srik} (Guangyun, in Baxter's transcription), which glosses the second as the colour of the face.
\end{cousinweb}

\begin{grammarlens}[title={What you've built}]
One more word for getting around.
\end{grammarlens}

\subsection*{Your turn}
\begin{itemize}
\raggedright
  \item \emph{Say it:} \emph{tèsè}
  \item \emph{Say it:} \emph{tèsè}, once more
  \item \emph{Say it:} \emph{zhè jiā fànguǎn de cài hěn yǒu tèsè} — The food at this restaurant has a character all its own.
  \item \emph{Recall:} say the Mandarin for a restaurant's signature dish, then the Mandarin for to eat no meat, then say \emph{tèsè} again
\end{itemize}

\begin{tcolorbox}[breakable,colback=teal!4,colframe=teal!35!black,title={Before you move on}]
What does \zh{特色} mean? (\textquotedblleft{}Special character; speciality\textquotedblright{}.) Say it once more.
\end{tcolorbox}
\section[\texorpdfstring{\zh{色} (sè)}{色 (sè)}]{\zh{色} — a small top piece, then a frame with a long hook}

\label{lesson:ZH-W451-se-colour}



\begin{quote}
The word first, then the character inside it.

\begin{itemize}
\raggedright
  \item \emph{Recall:} say \emph{tèsè} — \textbf{R1}, one lesson back, before it decays
\end{itemize}

\begin{quote}
\zh{包} \zh{特}
\end{quote}

\emph{Read it:} each one aloud, with nothing said first — \emph{bāo}, \emph{tè}
\end{quote}

\subsection*{Script — a small top piece, then a frame with a long hook}
\hlblockfigure{\detokenize{figures/ZH-W451-se-colour-filmstrip.pdf}}{How \zh{色} is written, stroke by stroke}

\begin{quote}
\zh{色}
\end{quote}

\textbf{Two pieces, top before bottom.} Count the pieces, not the strokes.

\textbf{The top piece, two strokes.} A short stroke falling to the left. Then a small bar that turns down to the left under it.

\textbf{The bottom piece, four strokes.} A bar that turns down on the right. A short vertical in the middle. A bar across under them. Last, a long stroke that runs down the left side, turns right along the foot and hooks up.

2 + 4 = six.

\subsection*{Your turn}
\begin{itemize}
\raggedright
  \item \emph{Write it:} \zh{色}, all six, with the model in front of you
  \item \emph{Cover:} write it again from memory, and count the pieces, not the strokes
\end{itemize}

\begin{tcolorbox}[breakable,colback=teal!4,colframe=teal!35!black,title={Before you move on}]
How many strokes in \zh{色}? (\textbf{Six}.) What is the last stroke? (\textbf{A long stroke down the left that turns along the foot and hooks up}.)

Source: \href{https://raw.githubusercontent.com/chanind/hanzi-writer-data/68d10a4b21150cae5e1ebbd223eed289cf32d90c/data/\%E8\%89\%B2.json}{Hanzi Writer Data \zh{色}, pinned stroke paths and medians}.
\end{tcolorbox}
\section[\texorpdfstring{\zh{出色} (chūsè)}{出色 (chūsè)}]{\zh{出色} (chūsè) — outstanding; excellent}

\label{lesson:ZH-C451-chuse}



\begin{quote}
Before the new one: say the Mandarin for to eat no meat, then the Mandarin for a special character of its own.
\end{quote}

\subsection*{You'll want to know: \zh{出色}}
\textbf{\zh{出色}} — \emph{chūsè} — \textquotedblleft{}outstanding; excellent\textquotedblright{}.

\begin{quote}
\textbf{\zh{他的中文很出色}} \emph{tā de Zhōngwén hěn chūsè} — His spoken Chinese is outstanding.
\end{quote}

\begin{cousinweb}
\textbf{\zh{出色}} (\emph{chūsè}) is built from \zh{出} 'to stand out' and \zh{色} 'a look', both inherited from Middle Chinese, \emph{tsyhwit} and \emph{srik} (Guangyun, in Baxter's transcription); how the pair came to mean 'outstanding' is not settled.
\end{cousinweb}

\begin{grammarlens}[title={What you've built}]
One more word for getting around.
\end{grammarlens}

\subsection*{Your turn}
\begin{itemize}
\raggedright
  \item \emph{Say it:} \emph{chūsè}
  \item \emph{Say it:} \emph{chūsè}, once more
  \item \emph{Say it:} \emph{tā de Zhōngwén hěn chūsè} — His spoken Chinese is outstanding.
  \item \emph{Recall:} say the Mandarin for to eat no meat, then the Mandarin for a special character of its own, then say \emph{chūsè} again
\end{itemize}

\begin{tcolorbox}[breakable,colback=teal!4,colframe=teal!35!black,title={Before you move on}]
What does \zh{出色} mean? (\textquotedblleft{}Outstanding; excellent\textquotedblright{}.) Say it once more.
\end{tcolorbox}
\section[sùcài]{sùcài — meatless dish; vegetable dish}

\label{lesson:ZH-C451-sucai}



\begin{quote}
Before the new one: say the Mandarin for a special character of its own, then the Mandarin for outstanding.
\end{quote}

\subsection*{You'll want to know: sùcài}
\emph{sùcài} — \textquotedblleft{}meatless dish; vegetable dish\textquotedblright{}.

\begin{quote}
\emph{zhè jiā fànguǎn yǒu hěn duō sùcài} — This restaurant has lots of meatless dishes.
\end{quote}

Its characters come later; for now it is a word for the ear.

\begin{cousinweb}
\emph{sùcài} is built from \emph{sù} 'plain' and \zh{菜} 'a dish', both inherited from Middle Chinese, \emph{suH} and \emph{tshojH} (Guangyun, in Baxter's transcription); the sense 'meatless' for the first is not given there.
\end{cousinweb}

\begin{grammarlens}[title={What you've built}]
One more word for getting around.
\end{grammarlens}

\subsection*{Your turn}
\begin{itemize}
\raggedright
  \item \emph{Say it:} \emph{sùcài}
  \item \emph{Say it:} \emph{sùcài}, once more
  \item \emph{Say it:} \emph{zhè jiā fànguǎn yǒu hěn duō sùcài} — This restaurant has lots of meatless dishes.
  \item \emph{Recall:} say the Mandarin for a special character of its own, then the Mandarin for outstanding, then say \emph{sùcài} again
\end{itemize}

\begin{tcolorbox}[breakable,colback=teal!4,colframe=teal!35!black,title={Before you move on}]
What does sùcài mean? (\textquotedblleft{}Meatless dish; vegetable dish\textquotedblright{}.) Say it once more.
\end{tcolorbox}
\section[shèng]{shèng — to be left over}

\label{lesson:ZH-C451-sheng}



\begin{quote}
Before the new one: say the Mandarin for outstanding, then the Mandarin for a meatless dish.
\end{quote}

\subsection*{You'll want to know: shèng}
\emph{shèng} — \textquotedblleft{}to be left over\textquotedblright{}.

\begin{quote}
\emph{shèng le hěn duō cài wǒmen dǎbāo} — There's a lot of food left, so we'll take it away.
\end{quote}

Its characters come later; for now it is a word for the ear.

\begin{cousinweb}
\emph{shèng} is inherited from Middle Chinese \emph{zyingH} (Guangyun, in Baxter's transcription), glossed as what is long and to spare.
\end{cousinweb}

\begin{grammarlens}[title={What you've built}]
One more word for getting around.
\end{grammarlens}

\subsection*{Your turn}
\begin{itemize}
\raggedright
  \item \emph{Say it:} \emph{shèng}
  \item \emph{Say it:} \emph{shèng}, once more
  \item \emph{Say it:} \emph{shèng le hěn duō cài wǒmen dǎbāo} — There's a lot of food left, so we'll take it away.
  \item \emph{Recall:} say the Mandarin for outstanding, then the Mandarin for a meatless dish, then say \emph{shèng} again
\end{itemize}

\begin{tcolorbox}[breakable,colback=teal!4,colframe=teal!35!black,title={Before you move on}]
What does shèng mean? (\textquotedblleft{}To be left over\textquotedblright{}.) Say it once more.
\end{tcolorbox}
